% Supporting proof fragment. This is deliberately not a standalone document.
\section{A positive-measure tensor transfer certificate}
\label{supp:positive-tensor-transfer}

I fix positive weights $w_i$ with $\sum_iw_i=1$, drift differences
$d_i$ and centred loadings $s_i\in\mathbb R^p$ satisfying
$\sum_iw_id_i=0$ and $\sum_iw_is_i=0$. I use the classically specified
risk-neutral model, and put
\[
 T(y)=\sum_iw_i e^{d_i+s_i\cdot y},\qquad U(y)=T(y)-1,
 \qquad m(y)=m_G+r\cdot y.
\]
Jensen's inequality gives $T\geq1$ and $U\geq0$ pointwise. For the
conditional payoff kernel $P$, define $C_2=P''-P'$ and
$C_{k+1}=C_k'-kC_k$ for $k\geq2$. I require finite certified bounds
\[
 M_k\geq\|C_k\|_\infty\ (k=2,3,4),\quad
 D_1\geq\|C_2'\|_\infty,\quad D_2\geq\|C_2''\|_\infty,
 \quad D_{31}\geq\|C_3'\|_\infty.
\]
These bounds exist for vanilla and capped lognormal call kernels, even
though the vanilla payoff itself is unbounded. They are payoff-specific
Gaussian boundary bounds, shared with all admitted pricing methods.

\subsection{One-coordinate transport as a positive measure}

Let $\nu(dy)=\phi(y)\mathbf1_{[-L,L]}(y)\,dy$ be the unnormalised
continuous retained Gaussian measure and let its mass be
$a=\Phi(L)-\Phi(-L)$. Divide its CDF range into $N$ equal-mass cells
$[a_c,b_c]$. Let $c_c$ be the actual dyadic rounded midpoint-quantile
node, with a certificate that $a_c\leq c_c\leq b_c$. Define the
unnormalised rounded measure
\[
 \tau=\frac aN\sum_c\delta_{c_c},\quad
 \mu_{1,c}=\int_{a_c}^{b_c}(y-c_c)\phi(y)\,dy,\quad
 \mu_{2,c}=\int_{a_c}^{b_c}(y-c_c)^2\phi(y)\,dy.
\]
Both $\nu$ and $\tau$ have exactly mass $a$. In particular $\tau$ is
not the normalised probability distribution over the quantum labels.

For $f\in C^2([-L,L])$, the fundamental theorem of calculus gives
\[
 f(y)-f(c)=f'(c)(y-c)+\int_c^y(y-z)f''(z)\,dz.
\]
For $y<c$, reversing the integral gives
$\int_y^c(z-y)f''(z)\,dz$, with a nonnegative coefficient. Fubini's
theorem is applicable because the integrands are continuous on a compact
cell. It gives
\[
 \int f\,d(\nu-\tau)
 =\sum_c\mu_{1,c}f'(c_c)+\int f''(z)\,\rho_2(dz),
\]
where $\rho_2$ is the finite positive measure with cellwise density
\[
 w_{a_c}(z)=\int_{a_c}^{z}(z-y)\phi(y)\,dy
 \quad(a_c\leq z\leq c_c),\qquad
 w_{b_c}(z)=\int_z^{b_c}(y-z)\phi(y)\,dy
 \quad(c_c\leq z\leq b_c).
\]
The expressions are
$w_e(z)=\phi(z)+z\Phi(z)-\phi(e)-z\Phi(e)$ on their respective
segments. They are nonnegative directly from their integral definitions.
Applying the identity to $f(y)=y^2/2$ proves
\[
 \rho_2([-L,L])=\frac12\sum_c\mu_{2,c}.
\]
Choose rational $\widehat\mu_{1,c}\geq|\mu_{1,c}|$ from outward
interval arithmetic and put
$\rho_1=\sum_c\widehat\mu_{1,c}\delta_{c_c}$. Then
\[
 \left|\int f\,d(\nu-\tau)\right|
 \leq\int|f'|\,d\rho_1+\int|f''|\,d\rho_2.
\]
The linear allowance intentionally discards cancellation. It never
attributes a favourable sign to an uncertain tiny first-moment interval.

For a finite measure $\eta$, write
$\mathcal M_\eta(v)=\int e^{vy}\,\eta(dy)$. Applying the exact transport
identity to $f(y)=e^{vy}$ gives, for $v\ne0$,
\begin{equation}
 \mathcal M_{\rho_2}(v)=
 \frac{\mathcal M_\nu(v)-\mathcal M_\tau(v)
       -v\sum_c\mu_{1,c}e^{vc_c}}{v^2}.
 \label{supp:eq:positive-transport-mgf}
\end{equation}
The sign follows from continuous minus rounded transport. Explicitly,
\[
 \mathcal M_\nu(v)=e^{v^2/2}\{\Phi(L-v)-\Phi(-L-v)\},\quad
 \mathcal M_\tau(v)=\frac aN\sum_c e^{vc_c},\quad
 \mathcal M_{\rho_1}(v)=\sum_c\widehat\mu_{1,c}e^{vc_c}.
\]
At $v=0$ I use $\mathcal M_{\rho_2}(0)=\sum_c\mu_{2,c}/2$
directly. An almost-zero exact rational loading can cause cancellation
in the quotient even when $v\ne0$. Positivity and support instead give
\[
 e^{-|v|L}\mathcal M_{\rho_2}(0)
 \leq\mathcal M_{\rho_2}(v)
 \leq e^{|v|L}\mathcal M_{\rho_2}(0).
\]
This fallback is rigorous, rather than a numerical substitution of zero.

\subsection{Tensor telescoping and derivative envelopes}

Product measures satisfy the exact algebraic identity
\[
 \nu^{\otimes p}-\tau^{\otimes p}
 =\sum_{j=1}^p\nu^{\otimes(j-1)}\otimes(\nu-\tau)
             \otimes\tau^{\otimes(p-j)}.
\]
For $k\in\{1,2\}$, define the finite positive hybrid measure
\[
 \Lambda_{j,k}=\nu^{\otimes(j-1)}\otimes\rho_k
                         \otimes\tau^{\otimes(p-j)},\qquad
 \langle h\rangle_{j,k}=\int h\,d\Lambda_{j,k}.
\]
Applying one-coordinate transport conditionally and using the triangle
inequality yields, for a twice continuously differentiable residual $R$,
\begin{equation}
 \left|\int R\,d\nu^{\otimes p}-\int R\,d\tau^{\otimes p}\right|
 \leq\sum_j\left\{\langle|\partial_jR|\rangle_{j,1}
                    +\langle|\partial_{jj}R|\rangle_{j,2}\right\}.
 \label{supp:eq:tensor-positive-transport}
\end{equation}
Compact support justifies every conditional integral and change of order.
No omitted factor of $a^{p-1}$ is permitted: it is already present in the
other-coordinate hybrid measures.

The second correction has the exact ordinary integral representation
\[
 R_2(y)=U(y)^2\int_0^1(1-z)t^{-2}
                 C_2(m(y)+\log t)\,dz,\qquad t=1+zU(y).
\]
For a fixed coordinate put $u=T_j$ and $v=T_{jj}\geq0$. For
$G=t^{-2}C_2(m+\log t)$, direct differentiation gives
\begin{align*}
 G_j&=r_jt^{-2}C_2'+zu t^{-3}C_3,\\
 G_{jj}&=r_j^2t^{-2}C_2''+2r_jzu t^{-3}C_3'
                       +zv t^{-3}C_3+z^2u^2t^{-4}C_4.
\end{align*}
The $C_3$ and $C_4$ terms use
$\partial_t(t^{-2}C_2(m+\log t))=t^{-3}C_3$ and
$\partial_t(t^{-3}C_3(m+\log t))=t^{-4}C_4$.
The product rule gives
\[
 \partial_j(U^2G)=2UuG+U^2G_j,
\quad
 \partial_{jj}(U^2G)=2(u^2+Uv)G+4UuG_j+U^2G_{jj}.
\]
Differentiation under the compact $z$ integral is valid. Since $t\geq1$,
the elementary integrals
$\int_0^1(1-z)dz=1/2$, $\int_0^1z(1-z)dz=1/6$ and
$\int_0^1z^2(1-z)dz=1/12$ prove the pointwise bounds
\begin{align*}
 |\partial_jR_2|
 &\leq M_2U|u|+\frac{|r_j|D_1}{2}U^2
                  +\frac{M_3}{6}U^2|u|,\\
 |\partial_{jj}R_2|
 &\leq M_2(u^2+Uv)+2|r_j|D_1U|u|
       +\frac{2M_3}{3}Uu^2+\frac{r_j^2D_2}{2}U^2\\
 &\hspace{5mm}+\frac{|r_j|D_{31}}{3}U^2|u|
       +\frac{M_3}{6}U^2v+\frac{M_4}{12}U^2u^2.
\end{align*}
In particular these allowances retain powers of the arithmetic to
geometric dispersion $U$; they do not replace it by its maximum on a
large Cartesian box.

Cauchy's inequality under any finite positive measure gives
\[
 \langle U|u|\rangle\leq
       \sqrt{\langle U^2\rangle\langle u^2\rangle},\qquad
 \langle U^2|u|\rangle\leq
 \min\!\left\{
   \sqrt{\langle U^2\rangle\langle U^2u^2\rangle},
   \sqrt{\langle U^4\rangle\langle u^2\rangle}
 \right\}.
\]
Neither inequality requires division by the total mass. I apply each to
the actual $\Lambda_{j,k}$ separately. Thus the transport allowance is
obtained from nonnegative polynomial moments
$U^2,u^2,Uv,Uu^2,U^2v,U^2u^2,U^4$, of degree at most four in the
asset exponentials.

\subsection{Exact factorisation, tails and computational scope}

Set $X_i=e^{d_i+s_i\cdot y}$. Any required polynomial is a finite sum
$h(y)=\sum_\alpha b_\alpha\prod_iX_i^{\alpha_i}$ with
$|\alpha|\leq4$. The coefficients $b_\alpha$ may have either sign.
Independence yields exactly
\[
 \langle h\rangle_{j,k}
 =\sum_\alpha b_\alpha e^{\sum_i\alpha_id_i}
      \mathcal M_{\rho_k}\!\left(\sum_i\alpha_i s_{ij}\right)
      \prod_{\ell<j}\mathcal M_\nu\!\left(\sum_i\alpha_i s_{i\ell}\right)
      \prod_{\ell>j}\mathcal M_\tau\!\left(\sum_i\alpha_i s_{i\ell}\right).
\]
I retain signed cross terms in $u^2$ and centred terms in $U^2$ rather
than applying absolute values to their polynomial coefficients.
Outward interval arithmetic encloses each sum. The established pointwise
nonnegativity permits intersecting its enclosure with $[0,\infty)$,
but does not permit replacing a negative upper endpoint by zero.

All signed or positive finite MGFs are evaluated by a finite exponential
series with an explicit total-variation remainder. At order $J$, with
$|c_c|\leq B$ and weights $b_c$, the remainder is bounded by
\[
 \left(\sum_c|b_c|\right)
          e^{|v|B}\frac{(|v|B)^{J+1}}{(J+1)!}.
\]
This follows termwise from Taylor's theorem for $e^x$ and the triangle
inequality. It includes signed centred first-moment weights and their
interval uncertainty. Near-zero-rate quotient fallbacks remain separate
from this series allowance.

For a standard independent Gaussian vector, Gaussian tilting gives
\[
 \mathbb E[e^{b+s\cdot Y}\mathbf1_{Y\notin[-L,L]^p}]
 =e^{b+\|s\|^2/2}
   \left[1-\prod_{j=1}^p\{\Phi(L-s_j)-\Phi(-L-s_j)\}\right].
\]
Expanding $U^2$ into degree-two exponential monomials therefore gives
an exact finite expression for $\mathbb E[U^2\mathbf1_{Y\notin[-L,L]^p}]$.
Since $|R_2|\leq M_2U^2/2$, its tail allowance is the product of that
moment and $M_2/2$. Combining it with
\eqref{supp:eq:tensor-positive-transport} bounds
\[
 \left|\mathbb ER_2(Y)-a^p\frac1{N^p}
                   \sum_{c_1,\ldots,c_p}R_2(c_1,\ldots,c_p)\right|.
\]
The price is absent from construction. The sum in this statement defines
the financial grid target; the proof evaluates its transfer allowance
without enumerating its $N^p$ points.

There are at most $\binom{n+4}{4}$ monomials of total degree at most four.
For fixed derivative order, cached scalar MGFs and polynomial sums replace
Cartesian enumeration by a number of monomial terms polynomial in the
asset count. This is a bound-construction statement. It gives neither a
dimension-free error constant nor a quantum pricing speed-up, and it
does not remove the cost of preparing or applying the quantum circuit.

\subsection{A matched first-correction certificate}

The first correction obeys $R_1=R_2+UP'(m)$. Differentiating the extra
term gives exactly
\[
 \partial_j(UP')=uP'+r_jUP'',\qquad
 \partial_{jj}(UP')=vP'+2r_juP''+r_j^2UP'''.
\]
For the cap I use certified global bounds $M_1\geq\|P'\|_\infty$,
$B_2\geq\|P''\|_\infty$ and $B_3\geq\|P'''\|_\infty$. Valid
fallbacks are $B_2=M_1+M_2$ and $B_3=M_1+3M_2+M_3$, because
$P''=P'+C_2$ and $P'''=P'+3C_2+C_3$. The additional allowances are
\[
 M_1\langle|u|\rangle_{j,1}+|r_j|B_2\langle U\rangle_{j,1},
\]
\[
 M_1\langle v\rangle_{j,2}+2|r_j|B_2\langle|u|\rangle_{j,2}
                   +r_j^2B_3\langle U\rangle_{j,2},
\]
where $\langle|u|\rangle\leq
\sqrt{\langle1\rangle\langle u^2\rangle}$.

For the vanilla call, put $G(y)=K e^{m(y)+\lambda^2/2}$. Its delta
satisfies $0\leq P'\leq G$, and the same identities give
$|P''|\leq G+M_2$ and $|P'''|\leq G+3M_2+M_3$. Thus the extra
first and second derivative bounds are respectively
\[
 G|u|+|r_j|GU+|r_j|M_2U,
\]
\[
 Gv+2|r_j|G|u|+2|r_j|M_2|u|
             +r_j^2GU+r_j^2(3M_2+M_3)U.
\]
The positive tilted measure $e^{m_G+\lambda^2/2+r\cdot y}
\Lambda_{j,k}(dy)$ allows the same factorised moments and Cauchy bounds.
Its MGFs use the exact shifted rates $r_\ell+\sum_i\alpha_is_{i\ell}$.

Both payoffs satisfy $0\leq R_1\leq GU$. The first-residual tail is
therefore bounded by the Gaussian tilted outside moment of $GU$;
for the cap I take its minimum with $H\mathbb P(Y\notin[-L,L]^p)$.
This tail replaces the second-residual tail. Adding both tails would
be valid but unnecessarily weak. No direct comparison may discard the
first-correction competitor solely because its transfer proof uses a
different residual order.

\subsection{Retained numerical scope}

At the frozen $L=6$, $b_y=10$, $q=28$, all 24 four- and eight-asset
capped and vanilla second-correction transfer bounds are below $0.005$.
The largest is approximately $0.002283696$ for the four-asset cap
with $\lambda=0.05$ and $H=25$. All 24 matched first-correction bounds
also pass that allocation; the largest is approximately $0.004913704$
for the same parameter case. The two-asset global-envelope construction
retains three allocation failures. Its stronger previously certified
local transport bounds remain available, so those failures are not
relabelled as missing two-asset financial prices. These are deterministic
financial transfer findings, distinct from complete algorithm budgets,
practical advantage, empirical validation and literature priority.
